\multlinegap0pt

\AtBeginDocument{%
\renewcommand\theenumi{\alph{enumi}}%
\renewcommand\theenumii{\roman{enumii}}%
}
\makeatletter%
\renewcommand\p@enumii{}%
\makeatother%

\let\epsilon\varepsilon
\newcommand\implique{\mathchoice{\Longrightarrow}{\Rightarrow}{\Rightarrow}{\Rightarrow}}
\newcommand{\spfrac}[2]{\sfrac{#1}{(#2)}}
\newcommand{\psfrac}[2]{\sfrac{(#1)}{#2}}
\newcommand{\pspfrac}[2]{\sfrac{(#1)}{(#2)}}
\newcommand{\Psfrac}[2]{(\sfrac{#1}{#2})}
\renewcommand{\eg}{e.g.,\xspace}
\newcommand{\bn}{\boldsymbol{n}}
\newenvironment{smallpmatrix}{\big(\begin{smallmatrix}}{\end{smallmatrix}\big)}
\renewenvironment{bmatrix}{\bigl(\begin{matrix}}{\end{matrix}\bigr)}

\newcommand{\llbracket}{[\![}
\newcommand{\rrbracket}{]\!]}

\usepackage{tikz}
\usetikzlibrary{arrows.meta, decorations.markings}
\definecolor{darkgreen}{rgb}{0,0.5,0}

\usepackage{array}

\usepackage{bezier}
\usepackage[centertags]{amsmath}

\usepackage{mathtools}
%\usepackage{stmaryrd}

\usepackage[mathcal]{euscript}

\theoremstyle{plain}
\newtheorem{theo}{Theorem}[section]
\newtheorem{lemm}[theo]{Lemma}
\newtheorem{coro}[theo]{Corollary}
\newtheorem{prop}[theo]{Proposition}
\theoremstyle{definition}
\newtheorem{defi}[theo]{Definition}
\newtheorem{remk}[theo]{Remark}
\newtheorem{expl}[theo]{Example}

\DeclareMathOperator{\Ker}{Ker}
\DeclareMathOperator{\rank}{rk}
\DeclareMathOperator{\range}{Ran}
\DeclareMathOperator{\id}{Id}

\newcommand{\suchthat}{\;|\:}
\newcommand{\midsuchthat}{\;\middle|\:}
\newcommand{\transp}{{\scriptscriptstyle\mathrm{T}}}

\newcommand{\norm}[1]{\left\lVert #1\right\lVert}
\newcommand{\abs}[1]{\left\lvert #1\right \rvert}
\newcommand{\floor}[1]{\left\lfloor #1 \right\rfloor}
\newcommand{\ceil}[1]{\left\lceil #1 \right\rceil}
\newcommand{\bnorm}[1]{\bigl\lVert #1\bigr\lVert}
\newcommand{\babs}[1]{\bigl\lvert #1\bigr \rvert}
\newcommand{\bfloor}[1]{\bigl\lfloor #1 \bigr\rfloor}
\newcommand{\bceil}[1]{\bigl\lceil #1 \bigr\rceil}
\newcommand{\Bnorm}[1]{\Bigl\lVert #1\Bigr\lVert}
\newcommand{\Babs}[1]{\Bigl\lvert #1\Bigr \rvert}
\newcommand{\Bfloor}[1]{\Bigl\lfloor #1 \Bigr\rfloor}
\newcommand{\Bceil}[1]{\Bigl\lceil #1 \Bigr\rceil}
\newcommand{\bgnorm}[1]{\biggl\lVert #1\biggr\lVert}
\newcommand{\bgabs}[1]{\biggl\lvert #1\biggr \rvert}
\newcommand{\bgfloor}[1]{\biggl\lfloor #1 \biggr\rfloor}
\newcommand{\bgceil}[1]{\biggl\lceil #1 \biggr\rceil}
\newcommand{\nnorm}[1]{\lVert #1\lVert}
\newcommand{\nabs}[1]{\lvert #1 \rvert}
\newcommand{\nfloor}[1]{\lfloor #1 \rfloor}
\newcommand{\nceil}[1]{\lceil #1 \rceil}

\newcommand{\scalprod}[2]{\left\langle #1, #2 \right\rangle}

\renewcommand\floatpagefraction{.9}
\renewcommand\topfraction{.9}
\renewcommand\bottomfraction{.9}
\renewcommand\textfraction{.1}
\setcounter{totalnumber}{50}
\setcounter{topnumber}{50}
\setcounter{bottomnumber}{50}

\newcommand{\Epsilon}{\mathrm E}

